%上下幅設定
\setlength{\textheight}{\paperheight}   % 本文領域を一旦紙面の高さにする
\setlength{\voffset}{-0.4truemm}      % 上の余白を39mm（=1inch+13.6mm）に
\setlength{\headheight}{0truemm}  % 
\setlength{\headsep}{0truemm}     % ヘッダの分だけ上の余白を小さくする
\addtolength{\textheight}{-1.8in} 

%左右幅設定
\setlength{\fullwidth}{\paperwidth}     % 全体の幅（テキスト領域の幅ではない）を一旦紙面幅にする
\setlength{\hoffset}{-0.4truemm}
\setlength{\evensidemargin}{0truemm} % 左の余白を25mm(=1inch-0.4mm)に
\setlength{\oddsidemargin}{0truemm}  % 上に同じ
\addtolength{\fullwidth}{-50truemm}     % 左右余白分だけ本文領域の幅を小さくする。
\setlength{\textwidth}{\fullwidth}      % テキスト幅を全体の幅に一致させる

% 数式,文
\renewcommand{\theequation}{\thesection.\arabic{equation}}
\renewcommand{\thefigure}{\thesection.\arabic{figure}}
\renewcommand{\thetable}{\thesection.\arabic{table}}
\makeatletter
\@addtoreset{equation}{section}
\@addtoreset{figure}{section}
\@addtoreset{table}{section}
\makeatother

\usepackage{amsmath,amsfonts,amssymb} 
\usepackage{amsthm} 
\usepackage[bbgreekl]{mathbbol}
  \DeclareSymbolFontAlphabet{\mathbb}{AMSb}
  \DeclareSymbolFontAlphabet{\mathbbl}{bbold}
  \DeclareMathSymbol{\bbepsilon}{\mathord}{bbold}{"0F}
\usepackage{bm}
\usepackage{mathtools}
\usepackage{physics}

\usepackage{centernot}

\usepackage{ytableau}
\ytableausetup{centertableaux,mathmode,boxsize=1em}

\allowdisplaybreaks


% 枠をかけるようにする
\usepackage{ascmac}

% 画像
\usepackage[dvipdfmx]{graphicx}
%\usepackage[format=hang,small,bf]{caption}
\usepackage[subrefformat=parens]{subcaption}
\captionsetup{compatibility=false}
\usepackage{ulem}
\usepackage{wrapfig}
\usepackage[pgf]{adjustbox}

%ハイパーリンク
\usepackage[dvipdfmx]{hyperref}
\usepackage[dvipdfmx]{pxjahyper}
\hypersetup{
  colorlinks,
  citecolor=TokiwaIro,
  linkcolor=RuriIro,
  urlcolor=RuriIro,
  linktocpage}
\setcounter{tocdepth}{3}

%cite
\usepackage{cite}


%===============================================================%
%color
%===============================================================%
\usepackage{color}
\definecolor{RuriIro}{rgb}{0.,0.28,0.60}
\definecolor{TokiwaIro}{rgb}{0.,0.39,0.16}
\definecolor{EnjiIro}{rgb}{0.70,0.25,0.29}

\definecolor{dred}{rgb}{0.7,0.15,0.09}
\definecolor{dblue}{rgb}{0,0.12,0.64}
\definecolor{dgreen}{rgb}{0.2,0.51,0.19}
\definecolor{pegn}{rgb}{0.33,0.51,0.14}

% 論文タイトル（subsection）の視認性を上げる
\usepackage{sectsty}
\subsectionfont{\color{RuriIro}}

\newcommand{\dred}[1]{{\color{dred} #1}}
\newcommand{\dblue}[1]{{\color{dblue} #1}}
\newcommand{\dgreen}[1]{{\color{dgreen} #1}}

\newcommand{\red}[1]{{\color{red} #1}}
\newcommand{\blue}[1]{{\color{blue} #1}}
\newcommand{\enji}[1]{{\color{EnjiIro} #1}}

\newcommand{\todo}[1]{{\color{EnjiIro} #1}}
%===============================================================%
%macro
%===========%
%数学記号%
%===============================================================%

\newcommand{\nn}{\nonumber}
%
\renewcommand{\mc}{\mathcal}
\newcommand{\mr}{\mathrm}
\newcommand{\ms}{\mathsf}
\newcommand{\mb}{\mathbf}
\newcommand{\mbb}{\mathbb}
%
% \newcommand*{\tr}[1]{\mathrm{tr}\,{#1}}
% \newcommand*{\Tr}[1]{\mathrm{Tr}\,{#1}}
\newcommand*{\re}[1]{\mathrm{Re}\,{#1}}
\newcommand*{\im}[1]{\mathrm{Im}\,{#1}}
\renewcommand{\Re}{\mathop{\mathrm{Re}}}
\renewcommand{\Im}{\mathop{\mathrm{Im}}}
\renewcommand*{\det}{\mathrm{det}\,}
\renewcommand*{\tr}{\mathrm{tr}\,}
%
\newcommand*{\del}{\partial}
\newcommand*{\diff}{\mr{d}}
% \newcommand*{\dd}{\mr{d}}
%
\newcommand*{\iu}{\mr{i}}
\newcommand*{\ee}{\mr{e}}
\renewcommand{\exp}[1]{\mathrm{exp}\left[ #1 \right]}
%
\newcommand*{\ol}[1]{\overline{#1}}
%
\renewcommand*{\vev}[1]{\left\langle{#1}\right\rangle}
% \newcommand*{\ket}[1]{\left|{#1}\right>}
% \newcommand*{\bra}[1]{\left<{#1}\right|}
\renewcommand{\abs}[1]{\left|{#1}\right|}
%

\newcommand{\SL}{\mathrm{SL}}
\newcommand{\GL}{\mathrm{GL}}
\newcommand{\SO}{\mathrm{SO}}
\newcommand{\SU}{\mathrm{SU}}
\newcommand{\rU}{\mathrm{U}}

\newcommand{\GeV}{\mathrm{GeV}}
\newcommand{\MeV}{\mathrm{MeV}}
\newcommand{\TeV}{\mathrm{TeV}}
\newcommand{\keV}{\mathrm{keV}}
\newcommand{\eV}{\mathrm{eV}}
\newcommand{\geV}{\mathrm{GeV}}
\newcommand{\meV}{\mathrm{MeV}}
\newcommand{\teV}{\mathrm{TeV}}
\renewcommand{\keV}{\mathrm{keV}}
\renewcommand{\eV}{\mathrm{eV}}
%Pfaffian%
\newcommand*{\pf}[1]{\text{Pf}\left({#1}\right)}
%===============================================================%
%ギリシャ数字%
%===============================================================%
\newcommand*{\one}{$\mathrm{I}$}
\newcommand*{\two}{$\mathrm{I}\hspace{-1.2pt}\mathrm{I}$}
\newcommand*{\three}{$\mathrm{I}\hspace{-1.2pt}\mathrm{I}\hspace{-1.2pt}\mathrm{I}$}
%===============================================================%
%その他%
%===============================================================%
\newcommand*{\kahler}{K\"{a}hler~}
\renewcommand*{\cp}{$\mathcal{CP}$}
\newcommand*{\mcp}{\mathcal{CP}}
\newcommand{\pl}{\mathrm{Pl}}
\newcommand{\vs}[1]{\vspace{{#1}pt}}
